\section{The exact derived comparison at every finite spectral thickening}
\label{DP:sec:dp-derived}

This calculation uses the balanced comparison of the original theta
complex, as constructed in the shared conversation at A1427, and the
polynomial-invertibility theorem proved in the accompanying balanced-kernel
chapter. It uses no replacement polynomial in place of an actual zeta
packet. Write $\mathscr O=\mathcal O_{\mathbb C,s_0}$ for the local ring of
holomorphic germs at the specified spectral coordinate $s_0$,
$A=\mathbb C[t]$, $M=\mathscr O\otimes_A Q$, and
\[
 Q=\mathscr B/\Theta V,\qquad g=2\xi,\qquad
 \beta:M\longrightarrow\mathscr O/(g),\qquad
 \beta(a\otimes[F])=[a\mathcal MF].
\]
The $A$-action is exactly $tF=DF=-x\partial_xF$ on the original
function space and $ta=sa$ on $\mathscr O$. Let $K=\ker\beta$.
The preceding calculation on these original spaces proves that
\[
 0\longrightarrow K\longrightarrow M\xrightarrow{\beta}
 \mathscr O/(g)\longrightarrow0
 \quad\hbox{is exact},\qquad
 (s-s_0):K\xrightarrow{\sim}K.                 \tag{DP.1}
\]
In particular this is not an assumption of injectivity of $\beta$;
the preceding chapter exhibits an explicit nonzero element of $K$.

\begin{proposition}[Derived finite spectral observations]
For every nonzero holomorphic germ $h\in\mathscr O$, the original
comparison induces a quasi-isomorphism
\[
 M\mathbin{\otimes^{\mathbf L}_{\mathscr O}}\mathscr O/(h)
 \xrightarrow{\ \beta\otimes1\ }
 (\mathscr O/(g))\mathbin{\otimes^{\mathbf L}_{\mathscr O}}
 \mathscr O/(h).                                      \tag{DP.2}
\]
Both the degree-minus-one kernel and the degree-zero quotient are
retained. The map commutes with multiplication by every holomorphic
germ, in particular the original action $a^s$ for every real $a>0$.
\end{proposition}
\begin{proof}
Write the specified germ as $h=(s-s_0)^nu$, where $n\geq0$ is its
actual vanishing order and $u(s_0)\ne0$. This identity retains the
entire unit $u$, not merely its value at $s_0$. Multiplication by $u$
is invertible on every $\mathscr O$-module. By (DP.1), multiplication
by $(s-s_0)^n$ is invertible on $K$. Thus the literal multiplication
map $h:K\to K$ is invertible.

The integral domain $\mathscr O$ has the free resolution
\[
 0\longrightarrow\mathscr O\xrightarrow{\times h}\mathscr O
 \longrightarrow\mathscr O/(h)\longrightarrow0.
\]
Place its two free terms in degrees $-1,0$. Tensoring it with the
three modules in (DP.1) gives a short exact sequence of complexes
\[
 0\longrightarrow[K\xrightarrow{h}K]
 \longrightarrow[M\xrightarrow{h}M]
 \longrightarrow[\mathscr O/(g)\xrightarrow{h}\mathscr O/(g)]
 \longrightarrow0.                                   \tag{DP.3}
\]
The first complex has zero cohomology in both degrees because its
differential is invertible. The cohomology exact sequence of (DP.3)
therefore proves (DP.2), including the degree $-1$ terms
$\ker(h:M\to M)$ and $\ker(h:\mathscr O/(g)\to\mathscr O/(g))$.
Its degree-zero terms are the actual quotients by $h$.
All maps in (DP.3) are multiplication or $\mathscr O$-linear maps;
they commute with multiplication by every germ. This proves the
claimed action compatibility without removing a nilpotent term.
\end{proof}

At an actual zero $\rho$ of complete order $m_\rho$, retain
$g(s)=(s-\rho)^{m_\rho}v_\rho(s)$ with its nonvanishing full unit.
For the exact thickening $h=(s-\rho)^n$, the right complex in
(DP.3) has
\[
 H^{-1}=\{z\in\mathcal O_\rho/(g):(s-\rho)^nz=0\},\qquad
 H^0=\mathcal O_\rho/(g,(s-\rho)^n).                 \tag{DP.4}
\]
When $n\geq m_\rho$, its differential is zero, so both terms are
the entire length-$m_\rho$ arithmetic local algebra. When
$0<n<m_\rho$, its kernel is exactly
$(s-\rho)^{m_\rho-n}\mathcal O_\rho/(g)$ and its quotient is
$\mathcal O_\rho/((s-\rho)^n)$, each of length $n$.
These statements follow by division by the retained unit
$v_\rho$, which is an isomorphism of the specified ideals and
does not change the action. On every retained full block the action
is still
\[
 a^{\rho}\sum_{j=0}^{m_\rho-1}
 \frac{(\log a)^j}{j!}N_\rho^j,\qquad
N_\rho[z^j]=[z^{j+1}],\quad z=s-\rho.                \tag{DP.5}
\]

Here is the exact bridge to the original dilation, rather than only
to scalar multiplication on a balanced module. Define
$R_aF(x)=F(x/a)$ and $R_a\phi(x)=\phi(x/a)$ for $a>0$.
These maps preserve the original spaces, commute with $D$, and
$\Theta R_a=R_a\Theta$ term by term. Substitution $x=ay$ gives
$\mathcal M(R_aF)(s)=a^s\mathcal MF(s)$; hence the induced
$\mathscr O$-linear map $1\otimes R_a$ on $M$ satisfies
$\beta(1\otimes R_a)=a^s\beta$. It preserves $M[h]$ and $hM$,
since it commutes with $h$. The proved comparison on kernels and
quotients therefore transports its action to literal multiplication
by $a^s$, giving (DP.5). This does not assert equality of
$1\otimes R_a$ and scalar $a^s$ on the entire balanced module:
their difference has image in $K$, which is retained.

For the counterfactual under examination, $\rho$ is an actual
off-critical nontrivial zero of $g$; its existence is the temporary
premise, not a conclusion. Equations (DP.2)--(DP.5) prove that its
complete finite block persists through this balanced comparison,
even after taking every finite derived spectral thickening. It
cannot be an element of the nonzero kernel $K$, since multiplication
by its annihilating polynomial is invertible on $K$.

On each of the three original nonempty two-leg masks whose admitted
theta image is the full $\Theta V$, these complexes have the split
lifts $(\lambda,x)\mapsto(\lambda,f(x))$, and external $\tau$ maps
to external $\tau$. A zero in (DP.3) is the supported zero at that
same mask. This assertion uses the full admitted source $V$. The
arithmetic spectral coordinate $s_0$ and the absolute base point
$\tau$ are connected by this labelled comparison; they are not
identified by it. Smaller admitted sources have the following
additional, explicitly computed finite spectral data.

\begin{proposition}[All finite defects of a smaller source label]
Let $W\subseteq V$ be an original $D$-invariant admitted source,
and retain $Q_W=\mathscr B/\Theta W$. Put
\[
 M_W=\mathscr O\otimes_AQ_W,\quad
 L_W=\mathscr O\otimes_A(V/W),\quad
 \beta_W(a\otimes[F]_W)=[a\mathcal MF],\quad K_W=\ker\beta_W.
\]
Then the source inclusion induces the exact sequences
\[
 0\longrightarrow L_W\xrightarrow{[\phi]\mapsto[\Theta\phi]}
 M_W\xrightarrow{\pi_W}M\longrightarrow0,\qquad
 0\longrightarrow L_W\longrightarrow K_W\longrightarrow K
 \longrightarrow0.                                    \tag{DP.6}
\]
For every nonzero germ $h$, their induced maps give
\[
 L_W[h]\xrightarrow{\sim}K_W[h],\qquad
 L_W/hL_W\xrightarrow{\sim}K_W/hK_W.                 \tag{DP.7}
\]
More fully, inclusion is a quasi-isomorphism
\[
 [L_W\xrightarrow{h}L_W]\xrightarrow{\sim}
 [K_W\xrightarrow{h}K_W].                              \tag{DP.8}
\]
\end{proposition}
\begin{proof}
Injectivity of the original $\Theta:V\to\mathscr B$ gives the
exact $A$-module sequence
$0\to V/W\xrightarrow{\Theta}Q_W\to Q\to0$. The flatness
proved in the balanced-kernel chapter yields the first sequence
in (DP.6). The equality $\beta_W=\beta\pi_W$ follows directly
from the retained Mellin formula. Every element of $L_W$ maps to
zero under $\beta_W$, since $\mathcal M\Theta\phi=gH_\phi$.
If $k\in K$, lift it to $x\in M_W$ using the surjection
$\pi_W$; then $\beta_Wx=\beta k=0$. This proves the second
sequence in (DP.6), including surjectivity of its last map.
Tensor its terms with the same two-term free resolution of
$\mathscr O/(h)$ used in (DP.3). The quotient complex is
$[K\xrightarrow{h}K]$, whose differential is invertible by
(DP.1) and its unit-factor argument. Its cohomology vanishes.
The resulting exact cohomology sequence proves both (DP.7) and
(DP.8). Every arrow comes from the actual source inclusion, and
commutes with the given actions.
\end{proof}

Thus a finite spectral class in a smaller label's Mellin kernel
is precisely a finite spectral class in its retained relation
quotient $V/W$. It vanishes upon passing to the full arithmetic
source $V$, by (DP.6). This is an actual transition map, not a
deletion of the smaller label or its extension. The next chapter
constructs such a nonzero class explicitly from the original
dilates of $\phi_*$, so that the restriction to the full source in
(DP.1)--(DP.5) is also tested by a concrete example. At a smaller
label only $D$-invariance has been imposed. Accordingly (DP.6)--(DP.8)
commute with the $D$ and holomorphic-scalar actions; an original
dilation $R_a$ acts between the labels $W$ and $R_aW$, and is an
endomorphism at $W$ only when $R_aW=W$. This follows from the exact
map $[F]_W\mapsto[R_aF]_{R_aW}$ and $\Theta R_a=R_a\Theta$.
